\section*{Resumen ejecutivo}

\textbf{El salario mínimo para 2027 debería aumentar en una proporción cercana a la inflación de 2026, sin superar el 7\% con las proyecciones disponibles.} El objetivo es compensar el aumento de los precios de este año sin añadir otra ganancia real de magnitud importante después del incremento de 2026. El equipo técnico del Banco de la República proyectó una inflación de cierre de 6,9\% en su informe de julio. La divulgación de la Encuesta de Opinión Financiera de octubre recoge una expectativa de 6,87\%. El porcentaje recomendado debe actualizarse con la información disponible al momento de la concertación.\cite{ipm,eof}

\textbf{Entre 2021 y 2025, el mínimo real aumentó 13,5\%, frente a un crecimiento de 1,2\% de la productividad por hora y una caída de 4,6\% por persona empleada.} El salario real histórico se calcula con el IPC promedio anual. La comparación acumula los cambios de 2022 a 2025 y no implica que el salario real haya subido todos los años. En 2022 perdió poder adquisitivo. Desde 2023 aumentó mientras la productividad laboral avanzó poco o retrocedió.\cite{ipc,productividad,salarios}

\textbf{El aumento nominal de 23\% en 2026 implica una ganancia real de aproximadamente 15,1\% si la inflación termina en 6,9\%.} El mínimo mensual pasó de \$1.423.500 a \$1.750.905, sin auxilio de transporte. La ganancia real de 2026 es todavía un escenario de cierre, calculado con IPC de diciembre. No es comparable directamente con una tasa de productividad de enero a septiembre ni demuestra cuál será la productividad anual de 2026.\cite{decreto2026,ipm}

\textbf{Limitar el ajuste al aumento de los precios permite conservar las ganancias reales anteriores sin repetir el salto de 2026.} Un incremento de 6,9\% llevaría el salario básico mensual a cerca de \$1.871.717. El techo de 7\% lo llevaría a \$1.873.468. Estas cifras excluyen el auxilio de transporte y no son una propuesta de congelamiento nominal. El 7\% es un límite provisional de política económica, no un porcentaje óptimo que pueda deducirse de la distancia entre dos índices.

\textbf{La protección de los trabajadores debe considerar también el acceso al empleo formal.} El mínimo aumenta la remuneración de quienes reciben el ajuste y conservan su empleo. Su efecto sobre quienes buscan trabajo depende de las respuestas de contratación, horas, precios y cumplimiento. La evidencia colombiana justifica atender estos riesgos, aunque la divergencia entre salario y productividad no permite medir por sí sola el efecto causal del aumento de 2026.\cite{espe,dube}
\clearpage
